% Generated draft: requires mathematical and source review.
\documentclass[11pt]{article}
\usepackage[margin=25mm]{geometry}
\usepackage{fontspec}
\setmainfont{Latin Modern Roman}
\usepackage{amsmath,amssymb,mathtools,unicode-math}
\setmathfont{Latin Modern Math}
\usepackage{enumitem,needspace,xurl,hyperref}
\setlength{\parindent}{0pt}
\setlength{\parskip}{5pt}
\setlength{\emergencystretch}{3em}
\newcommand{\N}{\mathbb{N}}
\newcommand{\Z}{\mathbb{Z}}
\newcommand{\Q}{\mathbb{Q}}
\newcommand{\R}{\mathbb{R}}
\newcommand{\C}{\mathbb{C}}
\newcommand{\F}{\mathbb{F}}
\newcommand{\A}{\mathbb{A}}
\newcommand{\PP}{\mathbb{P}}
\newcommand{\E}{\mathbb{E}}
\newcommand{\eps}{\varepsilon}
\newcommand{\dd}{\,\mathrm{d}}
\begin{document}
\section*{UnsolvedMath ID: 10600021}
\textbf{Source problem code:} catalogue entry
\subsection{Definitions and mathematical statement}
A long virtual knot is represented by an oriented long diagram agreeing with a straight oriented line outside a compact set, with classical and virtual crossings, modulo the corresponding Reidemeister and detour moves fixing the ends. Ordered concatenation defines an associative product \(\#\), with the straight long unknot as identity. A classical long knot admits a classical representative with no virtual crossings.

The source asks whether every commuting pair of long virtual knots has a common-power factorization up to classical factors:
\[
 K\#K'\simeq K'\#K
 \quad\Longrightarrow\quad
 \exists L,Q,Q',m,n:\qquad
 K\simeq L^m\#Q,\qquad K'\simeq L^n\#Q',
\]
where \(L\) is a long virtual knot, \(Q,Q'\) are classical long knots, and \(m,n\in\Z_{\geq0}\). Here \(\simeq\) is long virtual isotopy and \(L^m\) is the concatenation of \(m\) copies of \(L\), with \(L^0\) the identity. The same \(L\) must work for both knots. No uniqueness, primitivity or coprimality condition is imposed.

The source's motivation distinguishes the commutativity contributed by classical factors from general virtual noncommutativity. In the usual interpretation classical long factors commute with all long factors, while arbitrary virtual pairs need not commute. This background must not be read as saying that classical centrality itself fails. The proposed factorization is in the monoid, with actual classical factors, not only after passing to a group completion or quotient that could lose such factorization information.
\subsection{Short English statement}
If two long virtual knots commute under concatenation, must each be a power of one common long virtual knot followed by a classical long knot? Zero powers are allowed, and no uniqueness is requested.
\subsection{Sources}
\begin{itemize}
\item Local UnsolvedMath registry, permanent ID 10600021, source code catalogue entry; no proposer or proposal year supplied; updated 2026-08-06.
\item \emph{Unsolved Problems in Virtual Knot Theory and Combinatorial Knot Theory}, revised survey problem-list item 21: \url{https://arxiv.org/abs/1409.2823}. Registry access date 2026-07-29; source-TeX extraction. Provenance reports 60 top-level items, with compound items retained. Present primary-literature status and redistribution terms require review.
\item The statement refers to a paper by the Kamadas, Fenn et al. without an identifying citation. The summary lists ``Prime Decomposition and Non-Commutativity in the Monoid of Long Virtual Knots'' (2013), arXiv:1311.xxxx; that identifier is an incomplete placeholder, not a verified usable citation. No exact resolving reference is supplied.
\end{itemize}
\end{document}
